
%(BEGIN_QUESTION)
% Copyright 2009, Tony R. Kuphaldt, released under the Creative Commons Attribution License (v 1.0)
% This means you may do almost anything with this work of mine, so long as you give me proper credit

Beregn følgende verdier knyttet til oppvarming av en kjele med vann (2.1 pund aluminiumsgryte, 5.6 pund vann) fra 58 grader Fahrenheit til kokepunkt:

\vskip 10pt

\begin{itemize}
\item{} Varme som kreves for å nå koketemperatur = \underbar{\hskip 50pt} BTU
\vskip 5pt
\item{} Tid til koking (forutsatt 10,000 BTU/time varmeeffekt) = \underbar{\hskip 50pt} minutter
\vskip 5pt
\item{} Ekstra varme som kreves for å omdanne 2 pund vann til damp = \underbar{\hskip 50pt} BTU
\end{itemize}

\vskip 10pt

\noindent
Vis alt arbeid!

\vfil 

\underbar{file i00039}
\eject
%(END_QUESTION)





%(BEGIN_ANSWER)

Dette er et vurdert spørsmål -- ingen svar eller hint gis!

%(END_ANSWER)





%(BEGIN_NOTES)

Varme som kreves for å varme både gryta og vannet til koketemperatur kan beregnes med spesifikk varme-formelen:

$$Q = mc \Delta T$$

$$Q_{pot} = (2.1) (0.215) (212-58) = 69.531 \hbox{ BTU}$$

$$Q_{water} = (5.6) (1) (212-58) = 862.4 \hbox{ BTU}$$

$$Q_{total} = 69.531 + 862.4 = {\bf 931.9 \hbox{ BTU}}$$

\vskip 10pt

Med en varmeeffekt på 10,000 BTU per time (dvs. 166.67 BTU per minutt), vil dette ta {\bf 5.59 minutter} å nå koketemperatur.

\vskip 10pt

Når kokepunktet er nådd, kan varmen som kreves for å fordampe 2 pund vann beregnes med latent varme-formelen:

$$Q = mL$$

$$Q = (2)(970.3) = {\bf 1940.6 \hbox{ BTU}}$$


%INDEX% Physics, heat and temperature: calorimetry problem 

%(END_NOTES)
